\documentclass{article}
\usepackage{graphicx} % Required for inserting images

\title{location-service}
\author{Ahcène Amo}
\date{April 2026}

\begin{document}

\maketitle

\section{Introduction}

\section{Semaine 5 : Service de Localisation (Temps Réel)}

\subsection{Objectifs et Choix Technologiques}
Pour répondre aux exigences de haute fréquence (ingestion continue de dizaines de milliers de positions GPS), nous avons opéré une transition d'une architecture REST/CRUD classique vers une architecture \textit{Data-Intensive} (orientée données temps réel).
\begin{itemize}
    \item \textbf{Socle technique :} Node.js avec le framework NestJS (TypeScript) pour sa gestion optimale des entrées/sorties asynchrones (Event Loop).
    \item \textbf{Communication :} Utilisation du protocole gRPC (Protocol Buffers) en remplacement de HTTP/REST, afin de réduire considérablement l'empreinte réseau et la latence.
    \item \textbf{Base de données :} PostgreSQL couplé à l'extension TimescaleDB, indispensable pour le stockage performant des séries temporelles (Time-Series).
    \item \textbf{Bus d'événements :} Intégration asynchrone à Apache Kafka pour la propagation des états GPS (ex: franchissement de zones - Geofencing).
\end{itemize}

\subsection{Conception du Contrat gRPC}
L'interopérabilité entre les boîtiers matériels embarqués et notre serveur applicatif a été strictement définie via un contrat d'interface \texttt{location.proto}. Ce dernier expose un service \texttt{LocationService} comprenant :
\begin{itemize}
    \item \texttt{StreamPositions} : Une méthode de streaming bidirectionnel (\textit{Bidirectional Streaming}). Le véhicule pousse ses coordonnées en continu dans un tunnel persistant, et le serveur acquitte la réception des trames sans nécessiter la réouverture coûteuse d'une connexion TCP à chaque point.
    \item \texttt{WatchVehicle} : Une méthode de streaming serveur exploitant les \textit{Observables} (RxJS) de NestJS. Elle permet aux applications clientes (Tableaux de bord) de s'abonner en direct aux mouvements d'un véhicule spécifique.
\end{itemize}

\subsection{Persistance et Ingestion des Données}
Afin de supporter l'ingestion massive sans dégradation des performances en lecture, la table SQL \texttt{location\_readings} a été convertie en \textit{Hypertable} (spécificité TimescaleDB). Ce partitionnement automatique par plage temporelle garantit des requêtes historiques (API GraphQL) optimisées. En parallèle de la sauvegarde, chaque point valide est sérialisé et publié sur un topic Kafka pour alerter les autres microservices de la flotte.

\subsection{Simulation et Validation de l'Architecture}
Afin de valider la robustesse de la chaîne de bout en bout sans recourir à des véhicules physiques, un service indépendant de simulation a été conçu.
\begin{itemize}
    \item \textbf{Architecture du Simulateur :} Il s'agit d'un script Node.js autonome agissant comme un client gRPC externe. Cette séparation physique garantit que la validation s'effectue via le réseau et non en mémoire.
    \item \textbf{Comportement :} Le simulateur génère un itinéraire factice paramétré (ex: départ de Rouen) et émet des coordonnées GPS réalistes (incluant vitesse, cap, et précision) à intervalle régulier (3 secondes).
    \item \textbf{Déploiement Conteneurisé :} Intégré à la stack \texttt{docker-compose}, ce simulateur valide également la bonne résolution DNS interne (communication vers \texttt{location:50051}) et permet d'exécuter des tests de charge scalables sur la base de données temporelle.
\end{itemize}

\end{document}
